\documentclass{article}
\usepackage{iclr2027_conference,times}
\usepackage{amsmath,amssymb,amsthm,booktabs,graphicx,xcolor,hyperref,etoolbox}
\usepackage[T1]{fontenc}
\hypersetup{hidelinks}
\makeatletter
\patchcmd{\@maketitle}{Under review as a conference paper at ICLR 2027}{Research manuscript --- not submitted}{}{}
\patchcmd{\@maketitle}{Paper under double-blind review}{Draft prepared for independent review}{}{}
\makeatother
\newtheorem{theorem}{Theorem}
\newtheorem{proposition}[theorem]{Proposition}
\newtheorem{corollary}[theorem]{Corollary}
\newtheorem{lemma}[theorem]{Lemma}
\newcommand{\KL}{\operatorname{KL}}
\newcommand{\kl}{\operatorname{kl}}
\newcommand{\E}{\mathbb E}
\newcommand{\Prob}{\mathbb P}
\newcommand{\Alt}{\operatorname{Alt}}
\newcommand{\supp}{\operatorname{supp}}
\newcommand{\argmax}{\operatorname*{arg\,max}}
\newcommand{\argmin}{\operatorname*{arg\,min}}
\title{Grounding Only What Matters:\\Task-Relevant Permutation Cycles in\\Latent-Action World Models}
\author{Anonymous authors}
\date{}
\begin{document}
\maketitle
\lhead{Research manuscript --- not submitted}
\begin{abstract}
A world model can predict latent action effects without knowing which executable commands produce them. We study the additional interaction needed to ground only the effects required by a task. For a finite family of known latent transition kernels with an unknown shared command permutation, we characterize the optimal fixed-confidence information rate through directed permutation cycles that intersect the required effect set. This reduces experimental-design separation to a shortest-cycle problem and shows how task-required commands can be identified indirectly through the bijection. A likelihood confidence set supports task-sufficient stopping; a specialization of controlled sequential testing attains the information lower bound asymptotically. We construct multistep tasks for which full grounding costs $\Omega(\gamma^{-2}\log(1/\delta))$ while task grounding costs $O(\log(1/\delta))$, where $\gamma$ distinguishes irrelevant effects. We also derive a direct-sum law for independently aligned context components and a conservative model-error correction. Across 15,800 synthetic calibration runs, task-directed calibration separates sharply from full identification, but entropy sampling is competitive and faster at ordinary confidence. Estimated categorical kernels and deliberately inconsistent contexts reveal both invalid uncorrected certificates and substantial conservatism in robust certification. The results isolate the interaction cost of command grounding, conditional on useful latent effects and resettable diagnostic access, rather than solving representation discovery from arbitrary video.
\end{abstract}

\section{Introduction}
Learning a predictive effect and learning how to invoke it are different problems. A passive model may represent motion north, south, east, and west, yet not know the robot commands that realize those effects. Latent-policy and latent-action methods already address this distinction through a small amount of labeled interaction \citep{edwards2019,schmidt2023}. Recent identifiability results recover latent dynamics up to permutation under demonstrator diversity, and effect-alignment methods address inconsistent semantics across contexts \citep{schur2026,jiang2026}. The question here is not whether permutation ambiguity exists. It is \emph{which remaining ambiguities justify spending another calibration interaction}.

Full command identification can be unnecessarily demanding. Two visually distinguishable effects may be interchangeable for a task, or never be used by its latent policy. Conversely, information about an unused command can reveal a required command through the one-to-one mapping. Task relevance therefore does not justify simply discarding irrelevant actions from the calibration query space.

We formalize this tension in a restricted, auditable setting: finite latent effect kernels, a shared unknown permutation, resettable calibration contexts, and a known set of effects required by a downstream policy. Our main structural observation is that a competing command permutation is a collection of directed cycles. Only cycles intersecting the task-required effects can invalidate their grounding. This yields an exact task-dependent information-design problem with a shortest-cycle separation oracle.

Our statistical lower bound, likelihood-ratio confidence sets, and asymptotic adaptive testing argument specialize established controlled-sensing and pure-exploration tools \citep{nitinawarat2012,kaufmann2016,garivier2016}. The proposed contribution is the \emph{task-relevant permutation structure}, its consequences for experimental design and context alignment, and a controlled analysis of the gap between full and sufficient grounding. We do not claim that generic active hypothesis testing or latent-action calibration is new. Our experiments are mechanism tests, including strong baselines and failure cases, rather than evidence of state-of-the-art visual control.

\section{Conditional command grounding}
Let $\mathcal X$ be a finite set of calibration contexts, $\mathcal A=[m]$ the executable commands, and $P_j^x$ a distribution on a finite outcome alphabet $\mathcal Y$ for effect $j\in[m]$ in context $x$. Initially, the kernels are known, strictly positive, and share a common support. The unknown parameter $\theta\in\Theta=S_m$ is a bijection from commands to effects. A resettable query $q=(x,a)$ produces
\begin{equation}
Y\sim P_\theta^q:=P_{\theta(a)}^x. \label{eq:model}
\end{equation}
An adaptive algorithm chooses queries from prior observations and private randomness, stops at $\tau$, and outputs an answer. Calibration queries do not include the cost of navigating to $x$.

For a fixed nonempty required effect set $R$, define
\begin{equation}
d_R(\theta)=(\theta^{-1}(r))_{r\in R},\qquad
\Alt_R(\theta)=\{\lambda:d_R(\lambda)\ne d_R(\theta)\}. \label{eq:answer}
\end{equation}
A rule is $\delta$-correct if it stops almost surely and outputs $d_R(\theta)$ with probability at least $1-\delta$, for every $\theta$. Full grounding sets $R=[m]$. A capped implementation may instead abstain; its certificate remains valid, but the capped procedure is not an almost-sure identification algorithm.

\paragraph{Connection to control.}
Suppose a known latent policy uses only effects in $R$. Recovering $d_R$ makes that policy executable without identifying other commands. Our experimental deployment environment is an $H$-stage chain. At stage $h$, the required latent effect $r_h\in R$ advances with probability $p_g$; every other effect advances with $p_b<p_g$; failure terminates the episode with reward zero. Completing the chain gives reward one. The optimal value is $p_g^H$. Any deterministic policy wrong at a visited stage has regret at least
\begin{equation}
g_H=(p_g-p_b)p_g^{H-1}. \label{eq:gap}
\end{equation}
When the sequence covers $R$ and $\varepsilon<g_H$, grounding $d_R$ is equivalent to producing an $\varepsilon$-optimal deterministic chain policy. This equivalence is not asserted for arbitrary tasks with multiple acceptable policies; those require a more general answer formulation \citep{degenne2019}.

\paragraph{What is supplied.}
The latent kernels, required effect identities, task dynamics, and relative effect correspondences inside an aligned context component are inputs. A uniform passive command policy gives $m^{-1}\sum_jP_j^x$, independent of $\theta$, so additional data from that passive process alone cannot ground command labels. Recovering the effect kernels from general passive data is a separate identifiability problem \citep{schur2026,nikulin2025}. Section~\ref{sec:learned} examines estimated kernels under an explicit, strong anonymous-demonstrator assumption.

\section{Task-relevant information geometry}
For a distribution $w$ over allowed queries, let
\begin{align}
I_R(\theta,w)&=\min_{\lambda\in\Alt_R(\theta)}\sum_qw_q\KL(P_\theta^q\Vert P_\lambda^q),\\
D_R^*(\theta)&=\max_{w\in\Delta(\mathcal X\times\mathcal A)}I_R(\theta,w). \label{eq:rate}
\end{align}
The following standard change-of-measure specialization identifies the relevant statistical scale; a self-contained proof appears in Appendix~\ref{app:lower}.
\begin{theorem}[Information lower bound]\label{thm:lower}
For $0<\delta<1/2$, every $\delta$-correct rule with finite expected stopping time satisfies
\begin{equation}
\E_\theta\tau\ge\frac{\kl(1-\delta,\delta)}{D_R^*(\theta)}. \label{eq:lower}
\end{equation}
If a decision-changing alternative is observationally indistinguishable for every allowed query, no such finite-expectation rule exists.
\end{theorem}
The novelty question is therefore whether the command structure gives a sharper and more computable characterization of $D_R^*$ than generic hypothesis enumeration.

\subsection{Only cycles touching required effects matter}
Relabel a competing permutation as $\sigma=\lambda\circ\theta^{-1}$, acting on effect identities. Define the directed cost matrix
\begin{equation}
c_{ij}^{\theta}(w)=\sum_xw_{x,\theta^{-1}(i)}\KL(P_i^x\Vert P_j^x),\quad c_{ii}^{\theta}=0. \label{eq:cost}
\end{equation}
Let $\mathcal C_R$ contain simple directed cycles of length at least two whose vertices intersect $R$.
\begin{theorem}[Task-relevant cycle characterization]\label{thm:cycles}
For every $\theta$ and $w$,
\begin{equation}
I_R(\theta,w)=\min_{C\in\mathcal C_R}\sum_{(i,j)\in C}c_{ij}^{\theta}(w). \label{eq:cycles}
\end{equation}
Consequently $D_R^*$ is the optimum of a linear program with one inequality per relevant cycle. Given $w$, a most violated inequality can be found using all-pairs shortest paths in $O(m^3+|R|m)$ arithmetic operations after computing the costs.
\end{theorem}
\begin{proof}
The weighted divergence against $\lambda$ equals $\sum_i c_{i,\sigma(i)}^{\theta}$. The answer changes exactly when $\sigma$ moves some $r\in R$. Decompose $\sigma$ into disjoint cycles. Its cost is at least that of one nontrivial cycle touching $R$, because all costs are nonnegative. Conversely, any such cycle, with all other vertices fixed, is a decision-changing permutation. For the oracle, minimize $c_{rj}+\operatorname{dist}(j,r)$ over $r\in R,j\ne r$. Nonnegative costs permit a simple shortest return path.
\end{proof}
The oracle avoids enumerating $m!$ alternatives for design. Our implementation adds violated cycle constraints and checks it against the full alternative LP. This is a design result, not an end-to-end polynomial-time inference claim: the reference likelihood implementation still enumerates permutations.

\paragraph{Indirect grounding.}
The cost of a cycle depends on every command on that cycle, not only commands mapped into $R$. An optimal design can place zero mass on a required command while determining it by elimination. In our four-effect example at $\gamma=.02$, the identity-mapping task design places approximately $(0,.596,.175,.229)$ mass on commands at the strongest context. Effect $0$ is required but its own command need not receive positive asymptotic design mass. Forced exploration still samples it finitely and sublinearly.

\subsection{A sharp separation from complete grounding}
\begin{theorem}[Irrelevant-effect separation]\label{thm:separation}
There is a four-effect family with $R=\{0,1\}$ and an $H$-stage task, for any fixed $H\ge2$, such that for $0<\gamma\le .05$:
\begin{equation}
\E_\theta\tau_{\rm full}\ge
\frac{\kl(1-\delta,\delta)}{2\gamma\log((.5+\gamma)/(.5-\gamma))}, \label{eq:sepfull}
\end{equation}
whereas a task-grounding rule uses at most
$4\lceil32\log(16/\delta)\rceil$ queries. Thus full grounding has an $\Omega(\gamma^{-2}\log(1/\delta))$ requirement while task grounding admits an $O(\log(1/\delta))$ bound uniform in $\gamma$.
\end{theorem}
The strongest context returns two independent bits with success probabilities $(.8,.5)$, $(.2,.8)$, $(.2,.5-\gamma)$, and $(.2,.5+\gamma)$ for effects $0,1,2,3$. Distinguishing effects $2$ and $3$ requires resolving a vanishing gap. Identifying effects $0$ and $1$ uses fixed gaps. Appendix~\ref{app:separation} proves the constants and embeds these diagnostic queries and the multistep task in one model with a shared permutation. This is a separation between objectives, not a claim that every task-directed sampler beats every full-identification sampler at every budget.

\section{Calibrate until one answer is reliable}
Let $M=m!$, $\ell_t(\theta)=\sum_{s\le t}\log P_\theta^{q_s}(Y_s)$, and
\begin{equation}
B_\delta=\log((M-1)/\delta),\qquad
\mathcal H_t=\{\theta:\max_\lambda\ell_t(\lambda)-\ell_t(\theta)\le B_\delta\}. \label{eq:confidence}
\end{equation}
The set is recomputed at each time, not intersected across time. It is always nonempty.
\begin{proposition}[Anytime task certificate]\label{prop:certificate}
Under any adaptive query rule, $\Prob_\theta(\exists t:\theta\notin\mathcal H_t)\le\delta$. Therefore stopping once $d_R$ is constant on $\mathcal H_t$ and returning that answer has error probability at most $\delta$.
\end{proposition}
The proof applies Ville's inequality to each likelihood-ratio martingale and takes a union bound. No Bayesian credibility threshold is substituted for frequentist coverage. More generally, any common executable policy satisfying
\begin{equation}
\max_{\theta\in\mathcal H_t}(V_\theta^*-V_\theta^\pi)\le\varepsilon \label{eq:policy}
\end{equation}
is certified. It is essential to evaluate \emph{the same policy} under every surviving model, rather than checking each model's own optimal policy.

\paragraph{Task-directed query allocation.}
Fix a deterministic choice $w^*(h)$ of an optimal design in Equation~\eqref{eq:rate} for each candidate $h$. At time $t$, use the maximum-likelihood candidate $\hat\theta_{t-1}$ and draw
\begin{equation}
q_t\sim (1-t^{-1/2})w^*(\hat\theta_{t-1})+t^{-1/2}U, \label{eq:algorithm}
\end{equation}
where $U$ is uniform over queries. Update likelihoods and check Proposition~\ref{prop:certificate}. The known-kernel designs can be cached. We refer to this combination as TDG. The exploration term prevents permanent commitment to an initially wrong mapping.
\begin{theorem}[Asymptotic optimality, fixed identifiable instance]\label{thm:upper}
Assume all distinct candidate permutations are distinguishable by at least one allowed query and all kernels have a common finite strictly positive support. TDG stops almost surely, is $\delta$-correct, and
\begin{equation}
\lim_{\delta\downarrow0}\frac{\E_\theta\tau_\delta}{\log(1/\delta)}=\frac{1}{D_R^*(\theta)}. \label{eq:optimality}
\end{equation}
\end{theorem}
This is a task-answer specialization of controlled sequential testing, not a new generic optimality principle \citep{nitinawarat2012}. Appendix~\ref{app:upper} supplies an expectation-level proof: forced exploration gives summable maximum-likelihood error tails; thereafter the task-relevant log-likelihood ratios grow at the optimal rate. Constants depend on the full instance. In particular, Equation~\eqref{eq:optimality} is not uniform as irrelevant effects become identical; Theorem~\ref{thm:separation} provides the separate uniform-in-$\gamma$ result.

\subsection{Aligned context components}
Suppose known correct correspondences align contexts inside each of $c$ components, while the components have independent unknown command permutations. Every query belongs to one component, and the task requires its specified answer in every component.
\begin{corollary}[Direct-sum calibration cost]\label{cor:components}
Under the assumptions of Theorem~\ref{thm:upper} in each component, with local rates $D_k^*>0$, the optimal asymptotic total calibration cost is
\begin{equation}
\lim_{\delta\downarrow0}\frac{\inf_{\delta\text{-correct}}\E\tau_\delta}{\log(1/\delta)}
=\sum_{k=1}^c\frac1{D_k^*}. \label{eq:components}
\end{equation}
\end{corollary}
A lower bound changes only one component at a time. An upper bound runs local rules with error budgets $\delta/c$. This formalizes the value of reliable cross-context correspondences without asserting that visual similarity supplies such correspondences.

\section{What imperfect latent kernels permit}
Suppose nominal kernels $\widehat P$ obey, uniformly over candidates and outcomes,
\begin{equation}
|\log\widehat P_\theta^q(y)-\log P_\theta^q(y)|\le\eta_q,
\qquad E_t=\sum_{s\le t}\eta_{q_s}. \label{eq:logerror}
\end{equation}
\begin{proposition}[Robust calibration and control]\label{prop:robust}
Replace the likelihoods in Equation~\eqref{eq:confidence} by nominal likelihoods and its threshold by $B_\delta+2E_t$. The resulting set contains the true mapping at all times with probability at least $1-\delta$. If Equation~\eqref{eq:logerror} holds on a pretraining event of probability at least $1-\beta$, the joint guarantee is $1-\delta-\beta$. If nominal deployment kernels additionally have uniform total-variation error at most $\rho$, then a nominal common-policy regret bound of $\varepsilon-2H\rho$ certifies true regret at most $\varepsilon$ for horizon-$H$ terminal rewards in $[0,1]$.
\end{proposition}
For a fixed query allocation $w$, nominal expected log-likelihood separation is at least true separation minus $2\sum_qw_q\eta_q$. The robust threshold pays another $2\sum_qw_q\eta_q$. A positive quantity
\begin{equation}
I_R(\theta,w)-4\sum_qw_q\eta_q \label{eq:robustmargin}
\end{equation}
is a sufficient drift margin for eliminating task-changing alternatives under that fixed allocation. This is conservative, not an optimal robustness characterization. It explains why more calibration can fail to yield a certificate when the supplied latent model is insufficiently accurate. The experiments below use simulator-known radii to isolate this issue; they do not validate a deployable radius estimator.

\section{Experiments}\label{sec:experiments}
All experiments are synthetic, CPU-based, and evaluate chain values exactly. We record 15,800 calibration runs across five suites. This count includes multiple calibration runs sharing a pretrained model; independence is not asserted across those runs. Settings, seeds, caps, raw data, and reproduction commands are specified in Appendix~\ref{app:experiments} and the accompanying package.

\paragraph{Exact-model protocol.}
We use the four-effect construction above, $H=4$, $p_g=.8$, $p_b=.2$, and $R=\{0,1\}$. One context is diagnostic; eight others shrink both Bernoulli parameters toward $.5$ by a factor $.15$. Each run draws a uniformly random command permutation. We compare task and full KL designs under both task and full stopping, uniform sampling, and full-mapping posterior entropy reduction with task stopping. Nonuniform samplers share the forced-exploration schedule. Every method receives the same kernels and reset access, and paired random streams are used within a setting. Each setting has 200 replicas, $\delta=.05$, and a 50,000-query cap.

\begin{figure}[t]
\centering\includegraphics[width=\linewidth]{../figures/query_complexity.pdf}
\caption{Calibration cost versus separation of two unused effects. Error bars are 95\% bootstrap intervals over 200 replicas. The large separation from full recovery concerns a harder target; entropy/task is a strong baseline and is faster than TDG at this confidence level.}\label{fig:complexity}
\end{figure}

\begin{table}[t]
\centering\small% Automatically generated from raw experiment results. Do not edit numbers manually.
\begin{tabular}{lrrrr}
\toprule
Query design / stop & $\gamma=.01$ & $.02$ & $.04$ & $.08$ \\
\midrule
Task design / task stop & $78.2$ & $77.4$ & $81.6$ & $95.4$ \\
Full design / task stop & $625.8$ & $1031.0$ & $259.9$ & $123.2$ \\
Full design / full stop & $7817.0$ & $2258.2$ & $576.4$ & $179.5$ \\
Task design / full stop & $18839.6^{\dagger}$ & $4528.4$ & $1102.5$ & $202.4$ \\
Uniform / task stop & $590.0$ & $611.5$ & $653.3$ & $830.6$ \\
Entropy / task stop & $64.6$ & $68.6$ & $71.0$ & $92.3$ \\
\bottomrule
\end{tabular}

\caption{Mean calibration queries over 200 runs. Means include censoring at the cap; $\dagger$ denotes one censored run. All other cells have 200 certificates. ``Task design / task stop'' is TDG.}\label{tab:exact}
\end{table}

\paragraph{Full identification is expensive, but task allocation is not universally best.}
At $\gamma=.01$, TDG uses 78.2 queries on average (95\% bootstrap interval $[74.4,82.3]$), versus 7,817.0 for full/full ($[7,274.0,8,353.4]$). Changing only the stopping objective is helpful but incomplete: full/task uses 625.8. Changing only allocation gives task/full 18,839.6, with one cap hit. These ablations distinguish allocation from stopping. However, entropy/task uses only 64.6 queries, outperforming TDG. TDG produces one incorrect task certificate out of 200 in each nuisance setting; no zero-error claim is made. These observed fractions are below $.05$ but cannot independently validate an anytime guarantee.

\paragraph{Stringent confidence and the all-effects-required control.}
A separate two-context sweep at $\gamma=.02$ varies $\delta$. At $\delta=.1$, TDG and entropy/task use 64.4 and 55.4 queries; at $\delta=10^{-10}$ they use 224.0 and 375.5. This is consistent with a difference between finite-sample efficiency and an asymptotic design criterion, not proof of the asymptotic limit from finite data. When all four effects are required, task/full alternatives coincide: TDG and full/full are identical sample by sample, with mean 544.6 queries. Entropy/task uses 501.5. The negative control rules out attributing universal superiority to the task label alone.

\paragraph{Estimated kernels and misspecification.}\label{sec:learned}
For each of five source sample sizes, we generate 20 independent pretrained categorical models from anonymous pure demonstrators: each demonstrator consistently executes one effect, but its executable command label is absent. We estimate each effect kernel by smoothed counts, then run 24 calibration replicas per model, both with and without the oracle log-error correction. This is a deliberately favorable anchor assumption, not unsupervised video representation learning. At 4,096 source transitions per effect, robust calibration certifies only 36.0\% by 3,000 queries; at 65,536 it certifies all runs, averaging 90.6 queries. Uncorrected calibration averages approximately 61 queries but lacks the misspecified-model guarantee. Source sampling costs are reported separately and are not hidden inside the calibration count.

A second stress test mixes calibration kernels with a swap of effects 0 and 1 while leaving deployment unchanged. Strong mixing therefore creates inconsistent semantics across calibration and deployment, not a harmless global relabeling. At mixing weights $.7$ and $1$, uncorrected calibration certifies the wrong task controller in all 200 runs. The corrected rule gives no wrong certificates, but abstains in every nonzero-mixing condition by the cap, including mild shifts. Thus the current robustness result identifies a genuine safety--efficiency tradeoff and does not resolve practical model-error calibration.

\section{Related work and limitations}
ILPO and LAPO learn latent control representations and map them to executable actions with labeled interaction \citep{edwards2019,schmidt2023}. Identifiability from tagged demonstrators \citep{schur2026}, distractor-induced failures \citep{nikulin2025}, and cross-context effect alignment \citep{jiang2026} constrain when our conditional input model is plausible. We neither replace those representation methods nor claim that a few labels always suffice without their assumptions.

Our fixed-confidence arguments use controlled sensing, information-theoretic identification bounds, and time-uniform likelihood reasoning \citep{nitinawarat2012,kaufmann2016,garivier2016,howard2018}. Task-sufficient stopping is related to identification with multiple correct answers \citep{degenne2019}; our matching rate concerns the unique partial-permutation answer $d_R$, while Equation~\eqref{eq:policy} is only a general validity certificate. Novelty rests on the cycle restriction, structural implications, and grounded model-specific analysis, not these general statistical principles.

The present empirical scope is small finite stochastic systems and learned categorical likelihoods. There is no neural video model, physical robot, online navigation cost, learned task-relevance set, or validated data-dependent robust radius. The implementation enumerates candidate permutations; only the design oracle avoids factorial alternative enumeration. The task family is deliberately structured, so quantitative savings should not be extrapolated to general control interfaces. Independent human proof review and broader representation experiments remain necessary before treating this draft as submission-ready.

\section{Conclusion}
Minimal grounding is not necessarily complete command identification. In the finite shared-permutation model, its information constraints are exactly the directed cycles that touch task-required effects. This yields a precise account of indirect calibration, task-dependent stopping, and additive costs across independently aligned components. Controlled experiments support the separation from full recovery while exposing strong entropy baselines and the conservatism of robust certification. The next empirical step is to test the same cycle-informed calibration on imperfect latent effects learned from richer observations, without changing the information access or hiding representation failures.

\clearpage
\section*{Reproducibility and research status}
The accompanying package contains source code, deterministic seeds, all recorded calibration outcomes, summary tables, figure-generation code, and automated tests. Numerical results were produced by execution. Reported intervals are descriptive bootstrap intervals, not replacements for the analytic correctness guarantee. This is a research draft; its proofs have not yet received independent human review, and no acceptance, peer review, or conference submission is implied by its formatting. The public development repository must be replaced by an anonymized artifact before a double-blind submission.

\section*{AI use statement}
An OpenAI ChatGPT assistant was used for problem formulation, literature research, theorem and proof drafting, implementation, synthetic experimentation, analysis, and manuscript preparation. The experiments used programmatically sampled synthetic data, not invented numerical observations. Automated tests and an AI-assisted consistency review have been performed; independent human verification of mathematical correctness, novelty, experimental methodology, and citations has not yet been completed. The final human authors must verify and take responsibility for the complete submission and update this statement accordingly.

\section*{Ethics statement}
The experiments contain no human-subject data or physical deployment. A calibration certificate under a misspecified model should not be interpreted as a physical safety guarantee. The robust experiments explicitly expose conditions in which the method abstains rather than endorsing a controller.

\bibliographystyle{plainnat}
\bibliography{references}
\clearpage
\appendix
\section{Formal assumptions and notation}
The finite query set is $\mathcal Q=\mathcal X\times[m]$. At every time, the query is measurable with respect to past observations and fresh private randomness. Conditional on that query and history, the outcome has law $P_\theta^q$ and is independent of previous outcomes. All kernels have a common finite strictly positive support, so every log-likelihood increment is bounded in absolute value by a finite instance-dependent constant. A stopped transcript includes the queries, outcomes, stopping time, and output. The algorithm, including its randomization and stopping rule, is the same under every candidate parameter. Thus query-selection probabilities cancel in likelihood ratios. Natural logarithms are used throughout.

The required set $R$ is fixed before calibration. Its answer $d_R$ is a unique partial mapping, although multiple full parameters have that answer. This is a single-correct-answer identification problem at the level of the quotient answer space. It is not the most general multiple-correct-policy problem. The generic certificate in Equation~\eqref{eq:policy} remains valid for that broader problem, but the optimality theorem is not claimed for it.

\section{Information lower bound}\label{app:lower}
\begin{proof}[Proof of Theorem~\ref{thm:lower}]
Fix a true parameter $\theta$ and an alternative $\lambda\in\Alt_R(\theta)$. Write $N_q(\tau)$ for the query count. If $\E_\theta\tau=\infty$, the asserted lower bound is immediate. Otherwise, bounded log-likelihood increments and the stopped likelihood chain rule give
\begin{equation}
\KL(\mathsf P_\theta^{\rm transcript}\Vert\mathsf P_\lambda^{\rm transcript})
=\E_\theta\!\left[\ell_\tau(\theta)-\ell_\tau(\lambda)\right]
=\sum_q\E_\theta N_q(\tau)\KL(P_\theta^q\Vert P_\lambda^q). \label{eq:stoppedkl}
\end{equation}
For completeness, the final equality follows by writing the stopped sum as $\sum_{t\ge1}\mathbf1\{\tau\ge t\}\log(P_\theta^{q_t}(Y_t)/P_\lambda^{q_t}(Y_t))$, conditioning on the pre-outcome history and query, and exchanging expectation and sum. The absolute expectation is bounded by a constant times $\E_\theta\tau$.

Let $A$ be the event that the returned answer equals $d_R(\theta)$. Correctness gives $\mathsf P_\theta(A)\ge1-\delta$ and $\mathsf P_\lambda(A)\le\delta$. By data processing to $\mathbf1_A$, and monotonicity of binary relative entropy when the first argument exceeds the second,
\begin{equation}
\sum_q\E_\theta N_q(\tau)\KL(P_\theta^q\Vert P_\lambda^q)
\ge\kl(1-\delta,\delta).
\end{equation}
Normalize $w_q=\E_\theta N_q(\tau)/\E_\theta\tau$. The inequality holds for every task-changing alternative. Minimizing over these alternatives and upper-bounding the resulting rate by $D_R^*(\theta)$ proves Equation~\eqref{eq:lower}. If one such alternative has identical query laws, the left side is zero and contradicts $\kl(1-\delta,\delta)>0$.
\end{proof}
This is the standard stopped change-of-measure argument specialized to the partial-permutation answer; see \citet{kaufmann2016} for the general identification method.

\section{Cycle structure and computation}\label{app:cycles}
\begin{proof}[Detailed proof of Theorem~\ref{thm:cycles}]
For an alternative $\lambda$, set $\sigma(i)=\lambda(\theta^{-1}(i))$. Reindexing the commands by their true effects yields
\begin{align}
\sum_{x,a}w_{x,a}\KL(P_{\theta(a)}^x\Vert P_{\lambda(a)}^x)
&=\sum_i\sum_xw_{x,\theta^{-1}(i)}\KL(P_i^x\Vert P_{\sigma(i)}^x)\\
&=\sum_i c_{i,\sigma(i)}^\theta(w).
\end{align}
For $r\in R$, equality $\lambda^{-1}(r)=\theta^{-1}(r)$ is equivalent to $\lambda(\theta^{-1}(r))=r$, hence to $\sigma(r)=r$. Therefore $\lambda\in\Alt_R(\theta)$ exactly when the cycle decomposition of $\sigma$ contains a nontrivial cycle touching $R$.

All edge costs are nonnegative. Removing every cycle except one relevant cycle cannot increase total cost. Conversely, extend any relevant simple cycle by fixing all other vertices. The resulting permutation belongs to the alternative set and has exactly that cycle's cost. This proves equality of the two minima, including zero-cost or tied cases.

Introduce a variable $z$ for the worst-alternative information rate. For each relevant cycle $C$, impose
\begin{equation}
z\le\sum_{(i,j)\in C}\sum_xw_{x,\theta^{-1}(i)}\KL(P_i^x\Vert P_j^x),\qquad w\in\Delta(\mathcal Q).
\end{equation}
Maximizing $z$ is the desired LP. Given a tentative $w$, all-pairs shortest-path distances in the nonnegative complete directed cost graph give the minimum relevant nontrivial cycle as
\begin{equation}
\min_{r\in R}\min_{j\ne r}\{c_{rj}+\operatorname{dist}(j,r)\}.
\end{equation}
A shortest return path can be chosen simple: deleting a repeated-vertex subpath does not increase cost. It cannot visit $r$ before its endpoint. Thus prepending $(r,j)$ gives a valid simple relevant cycle. Computing distances costs $O(m^3)$ and inspecting the candidate roots costs $O(|R|m)$. Costs themselves require $O(|\mathcal X|m^2|\mathcal Y|)$ arithmetic operations if divergences are not precomputed.

Our cutting-plane implementation starts from relevant two-cycles, repeatedly solves the current LP, and adds a violated shortest cycle. It terminates after finitely many distinct cuts in exact arithmetic. We do not claim a polynomial bound on this particular cutting-plane iteration count. The polynomial-time separation oracle is a structural property; generic LP oracle reductions require their usual precision assumptions. Floating-point tests compare attained objectives, not only solver-reported objective values.
\end{proof}

\begin{corollary}[Task monotonicity]
If $R_1\subseteq R_2$, then $D_{R_1}^*(\theta)\ge D_{R_2}^*(\theta)$. Equality holds when their alternative sets coincide, including $R_1=R_2=[m]$.
\end{corollary}
\begin{proof}
Every cycle touching $R_1$ also touches $R_2$. The minimization has more constraints for $R_2$, so its optimal value cannot increase.
\end{proof}

\section{Full versus task grounding}\label{app:separation}
\begin{proof}[Proof of Theorem~\ref{thm:separation}]
At a diagnostic context, let $P_j$ be the product of two Bernoulli laws with parameter pairs
\begin{equation}
(.8,.5),\quad(.2,.8),\quad(.2,.5-\gamma),\quad(.2,.5+\gamma).
\end{equation}
They all have a common strictly positive support. Other allowed calibration contexts may independently pass each bit through a binary noise channel, replacing a parameter $p$ by $.5+s(p-.5)$ with $0\le s\le1$. These channels can only decrease relative entropy.

For any true $\theta$, form $\lambda$ by swapping the effects assigned to its commands for effects $2$ and $3$. This changes the full mapping and leaves $d_{\{0,1\}}$ unchanged. Queries to other commands have zero relative entropy. Queries to either swapped command at the strongest context have relative entropy
\begin{align}
k_\gamma
&=\KL(\operatorname{Ber}(.5-\gamma)\Vert\operatorname{Ber}(.5+\gamma))\\
&=(.5-\gamma)\log\frac{.5-\gamma}{.5+\gamma}
 +(.5+\gamma)\log\frac{.5+\gamma}{.5-\gamma}\\
&=2\gamma\log\frac{.5+\gamma}{.5-\gamma}.
\end{align}
The reverse divergence is equal by symmetry. Every allowed query therefore has divergence at most $k_\gamma$ for this pair of full mappings. Equation~\eqref{eq:stoppedkl} proves $\E\tau_{\rm full}\ge\kl(1-\delta,\delta)/k_\gamma$. A Taylor expansion gives $k_\gamma=8\gamma^2+O(\gamma^4)$.

For task grounding, sample each of the four commands $n=\lceil32\log(16/\delta)\rceil$ times at the strongest context. Estimate both bit means for each command. Hoeffding's inequality and a union bound over eight means give
\begin{equation}
\Prob\left(\max_{a,b}|\widehat p_{a,b}-p_{\theta(a),b}|\ge1/8\right)
\le16e^{-n/32}\le\delta.
\end{equation}
On the complementary event, effect $0$ is the unique command with first-bit mean above $.5$. For $\gamma\le .05$, effect $1$ is the unique command with second-bit mean above $.675$, since its true parameter is $.8$ while every other second-bit parameter is at most $.55$. When no unique threshold assignment exists, return any fixed answer; that possibility is already covered by the failure event. Thus $4n$ queries suffice uniformly in $\gamma$.

To embed control, adjoin a chain with states $0,\ldots,H$, an absorbing failure state, and required effects $0,1,0,1,\ldots$. The same unknown permutation determines which command realizes every latent effect in diagnostic and deployment contexts. Each correct effect advances with $p_g$ and every other effect with $p_b$. Calibration access is restricted to the diagnostic contexts. The model includes both diagnostic outcomes and chain dynamics, but the latter are not extra free calibration observations. Any deterministic policy misgrounding a required stage has value at most $p_b p_g^{H-1}$, proving the gap in Equation~\eqref{eq:gap}. Since both required effects occur when $H\ge2$, the decoded task answer is exactly the information needed to implement the optimal chain policy.
\end{proof}

\paragraph{Noiseless check.}
If each diagnostic outcome exactly reveals its effect, $m-1$ distinct command queries identify the full permutation, and $m-1$ are necessary in the worst case: two unqueried commands can still be exchanged. With $c$ independent components and exact known correspondences inside each, the analogous full-identification count is $c(m-1)$. These are elementary checks, not noisy finite-sample results or novelty claims. Even a single required effect can take $m-1$ noiseless queries in the worst case; task relevance does not always reduce worst-case counts.

\section{Anytime confidence and task certificates}\label{app:certificate}
\begin{proof}[Proof of Proposition~\ref{prop:certificate}]
For a false candidate $\lambda$, let
\begin{equation}
L_t^\lambda=\exp(\ell_t(\lambda)-\ell_t(\theta)).
\end{equation}
Conditioning on history and the selected query, the conditional expectation of its likelihood-ratio multiplier is
$\sum_yP_\theta^q(y)P_\lambda^q(y)/P_\theta^q(y)=1$. Hence $L_t^\lambda$ is a nonnegative mean-one martingale under the true model, including adaptive and randomized query selection. Ville's inequality gives
$\Prob_\theta(\sup_tL_t^\lambda>e^{B_\delta})\le e^{-B_\delta}$.

If $\theta\notin\mathcal H_t$, some $\lambda\ne\theta$ has $L_t^\lambda>e^{B_\delta}$. A union bound over the $M-1$ false candidates proves the uniform coverage claim. On that coverage event, every answer certified as constant on $\mathcal H_t$ must equal the true answer. The same event validates any common policy satisfying Equation~\eqref{eq:policy}, even when that policy is selected adaptively from the data. No termination claim is needed for the validity of a certificate; termination for TDG follows from the next theorem.
\end{proof}
Time-uniform supermartingale reasoning is standard; see \citet{howard2018}. The constant threshold is possible here because the hypothesis family is finite and every likelihood is fully specified. It should not be reused without correction when effect kernels are estimated from the calibration data itself.

\section{Expected-time asymptotic optimality}\label{app:upper}
\begin{proof}[Proof of Theorem~\ref{thm:upper}]
We analyze the infinite query process that follows Equation~\eqref{eq:algorithm} even after the potential stopping time. Fix $\theta$. For $\lambda\ne\theta$, define the per-query Hellinger deficit
\begin{equation}
h_{\theta\lambda}(q)=1-\sum_y\sqrt{P_\theta^q(y)P_\lambda^q(y)},\qquad
\bar h_{\theta\lambda}=|\mathcal Q|^{-1}\sum_qh_{\theta\lambda}(q).
\end{equation}
Full pairwise distinguishability implies $\eta:=\min_{\lambda\ne\theta}\bar h_{\theta\lambda}>0$. Conditional on the history, the query law assigns at least $t^{-1/2}/|\mathcal Q|$ mass to every query. Therefore
\begin{align}
\E_\theta\left[\sqrt{L_t^\lambda}\mid\mathcal F_{t-1}\right]
&\le\sqrt{L_{t-1}^\lambda}(1-t^{-1/2}\bar h_{\theta\lambda}),\\
\E_\theta\sqrt{L_t^\lambda}
&\le\exp\left(-\bar h_{\theta\lambda}\sum_{s=1}^t s^{-1/2}\right).
\end{align}
If the maximum-likelihood estimate is not $\theta$, at least one false likelihood is as large as the true one. Markov's inequality and a union bound yield constants $A,a>0$, depending on the fixed instance but not on $\delta$, such that
\begin{equation}
\Prob_\theta(\hat\theta_t\ne\theta)\le A e^{-a\sqrt t}. \label{eq:mletail}
\end{equation}
Thus errors are summable and the MLE is eventually correct almost surely. In particular, for an integer $u\ge1$,
\begin{equation}
\Prob_\theta(\exists s\ge u:\hat\theta_s\ne\theta)
\le A'(\sqrt u+1)e^{-a'\sqrt u} \label{eq:latemle}
\end{equation}
for fixed positive constants $A',a'$.

For $\lambda\in\Alt_R(\theta)$, write
\begin{equation}
Z_t^\lambda=\ell_t(\theta)-\ell_t(\lambda)
=\sum_{s=1}^t\mu_s^\lambda+M_t^\lambda,
\end{equation}
where $\mu_s^\lambda=\E_\theta[Z_s^\lambda-Z_{s-1}^\lambda\mid\mathcal F_{s-1}]$ and $M_t^\lambda$ is a martingale. Finiteness and positive support provide $K<\infty$ with every log-likelihood increment bounded by $K$ in absolute value. The martingale increments are bounded by $2K$. All predictable means are nonnegative. Whenever $\hat\theta_{s-1}=\theta$,
\begin{equation}
\mu_s^\lambda\ge(1-s^{-1/2})\sum_qw_q^*(\theta)\KL(P_\theta^q\Vert P_\lambda^q)
\ge(1-s^{-1/2})D_R^*(\theta).
\end{equation}
Let $D=D_R^*(\theta)>0$. On the event that all MLEs from time $u$ onward are correct,
\begin{equation}
\sum_{s=1}^t\mu_s^\lambda\ge D(t-u-2\sqrt t-2),\qquad t>u. \label{eq:drift}
\end{equation}

Fix $\zeta>0$ and set $\alpha=\zeta/[4(1+\zeta)]$, $u=\lceil\alpha t\rceil$. For all sufficiently large $t$ with $t\ge(1+\zeta)B_\delta/D$, Equation~\eqref{eq:drift} exceeds $B_\delta+\kappa t$ for some fixed $\kappa>0$ depending on $D,\zeta$ but not on $\delta$. For example, once the square-root and integer terms are absorbed, one may take $\kappa=D\zeta/[2(1+\zeta)]$.

If $Z_t^\lambda>B_\delta$ for every task-changing alternative, all those alternatives lie outside $\mathcal H_t$ because $\max_h\ell_t(h)\ge\ell_t(\theta)$. The set is nonempty and has a single task answer, so the stopping time is at most $t$. Combining Equation~\eqref{eq:latemle} and Azuma's inequality gives, in the stated range,
\begin{equation}
\Prob_\theta(\tau_\delta>t)
\le A''(\sqrt t+1)e^{-a''\sqrt t}
+|\Alt_R(\theta)|\exp\left(-\frac{\kappa^2t}{8K^2}\right). \label{eq:stoppingtail}
\end{equation}
Both sequences are summable, which establishes finite expectation and almost-sure termination. More precisely, summing the tail after $\lceil(1+\zeta)B_\delta/D\rceil$ yields an $o(B_\delta)$ contribution as $\delta\downarrow0$. Hence
\begin{equation}
\limsup_{\delta\downarrow0}\frac{\E_\theta\tau_\delta}{\log(1/\delta)}\le\frac{1+\zeta}{D}.
\end{equation}
Let $\zeta\downarrow0$. Theorem~\ref{thm:lower}, together with $\kl(1-\delta,\delta)/\log(1/\delta)\to1$, gives the matching lower limit. Proposition~\ref{prop:certificate} supplies $\delta$-correctness.
\end{proof}
The argument follows the architecture of controlled sequential hypothesis testing \citep{nitinawarat2012}. The full-identifiability assumption is stronger than necessary for task correctness and can be inefficient in nearly indistinguishable nuisance directions. It is used to obtain eventual full-MLE stabilization for this simple allocation rule, not because the task certificate needs full identification.

\subsection{A finite-time fixed-design companion}
\begin{proposition}[Task-only Hellinger upper bound]\label{prop:finite}
Draw calibration queries independently from a fixed $w$. Let
\begin{equation}
h_R(\theta,w)=\min_{\lambda\in\Alt_R(\theta)}\sum_qw_qh_{\theta\lambda}(q)>0,
\qquad K_R=|\Alt_R(\theta)|.
\end{equation}
Use the same task certificate. For every integer $n$,
\begin{equation}
\Prob_\theta(\tau>n)\le\min\{1,K_R\exp(B_\delta/2-nh_R(\theta,w))\}.
\end{equation}
In particular, with $n_0=\lceil(B_\delta/2+\log K_R)/h_R\rceil$,
$\E_\theta\tau\le n_0+(1-e^{-h_R})^{-1}$.
\end{proposition}
\begin{proof}
If $\tau>n$, some task-changing alternative has $Z_n^\lambda\le B_\delta$. Markov's inequality applied to $e^{-Z_n^\lambda/2}$ bounds this probability by $e^{B_\delta/2}(1-\sum_qw_qh_{\theta\lambda}(q))^n\le e^{B_\delta/2-nh_R}$. Union-bound over alternatives and sum the resulting geometric tail. Unlike Theorem~\ref{thm:upper}, this proposition needs no distinguishability between parameters with the same task answer.
\end{proof}

\section{Context components}\label{app:components}
\begin{proof}[Proof of Corollary~\ref{cor:components}]
The parameter is $\boldsymbol\theta=(\theta_1,\ldots,\theta_c)$ and the answer is the tuple of local task answers. A query in component $k$ depends only on $\theta_k$. Fix $k$ and compare the truth to alternatives that change only that component's answer. The stopped KL identity then contains only queries in component $k$. Normalize the expected counts in that component to obtain
\begin{equation}
\E_{\boldsymbol\theta}\tau_k\ge\frac{\kl(1-\delta,\delta)}{D_k^*}.
\end{equation}
Summing over components gives the lower bound, since $\tau=\sum_k\tau_k$. For an upper bound, run the local TDG procedures, each with error budget $\delta/c$, in any sequential order. A union bound gives joint correctness. Theorem~\ref{thm:upper} and fixed $c$ imply
\begin{equation}
\frac{\sum_k\E\tau_k(\delta/c)}{\log(1/\delta)}\longrightarrow\sum_k\frac1{D_k^*}.
\end{equation}
Known relative correspondences are needed to define a common local effect coordinate system. No result in this proof estimates those correspondences or their graph from observations.
\end{proof}

\section{Robustness to nominal likelihood error}\label{app:robust}
\begin{proof}[Proof of Proposition~\ref{prop:robust}]
Condition on a pretrained model and a valid uniform log-error bound. For every candidate,
$|\widehat\ell_t(\theta)-\ell_t(\theta)|\le E_t$. Therefore, if a false candidate exceeds the true nominal log-likelihood by more than $B_\delta+2E_t$, it exceeds the true exact log-likelihood by more than $B_\delta$. Excluding the true mapping from the robust confidence set implies one of the exact likelihood martingales crossed $e^{B_\delta}$. The same union bound as Proposition~\ref{prop:certificate} gives probability at most $\delta$. Conditioning on an independent pretraining data set for which the uniform bound holds with probability at least $1-\beta$, and adding its failure probability, gives $\delta+\beta$.

For the deployment statement, couple true and nominal transitions using maximal couplings at each step. Uniform total variation at most $\rho$ bounds the probability that the two trajectories ever disagree by $H\rho$ (or one if this exceeds one). Since the terminal reward is in $[0,1]$, every fixed policy's value differs by at most $H\rho$. The optimal values differ by the same amount by taking a supremum over the common executable policy class. For a selected policy $\pi$,
\begin{equation}
V_\theta^*-V_\theta^\pi
\le\widehat V_\theta^*-\widehat V_\theta^\pi+2H\rho.
\end{equation}
This inequality is pointwise in every policy, so it remains valid for a data-dependent policy on the robust coverage event. A nominal common-policy regret at most $\varepsilon-2H\rho$ completes the proof. When that tolerance is negative, this sufficient certificate cannot be issued.
\end{proof}

\paragraph{Why the correction can prevent stopping.}
For a fixed allocation $w$ and alternative $\lambda$, the exact expected log-likelihood ratio per query is $\sum_qw_q\KL(P_\theta^q\Vert P_\lambda^q)$. Replacing both terms in the ratio by nominal ones lowers this mean by at most $2\bar\eta_w$, where $\bar\eta_w=\sum_qw_q\eta_q$. The threshold grows at rate $2\bar\eta_w$. The strong law for i.i.d. query--outcome pairs therefore provides a sufficient positive drift of $I_R(\theta,w)-4\bar\eta_w$. If it is positive, each task-changing alternative is eventually excluded almost surely under that fixed design. A nonpositive lower bound is not an impossibility theorem; the correction can be loose. The robust experiments deliberately retain this conservative correction instead of interpreting failure to stop as successful identification.

\section{Experimental specification}\label{app:experiments}
\subsection{Environments, baselines, and accounting}
The calibration alphabet consists of four outcomes, representing two independent binary observations. With $p,z$ their Bernoulli parameters, the categorical probabilities are $(1-p)(1-z),(1-p)z,p(1-z),pz$. At the strongest context, the four parameter pairs are specified in Appendix~\ref{app:separation}. Weak contexts apply $p\mapsto.5+.15(p-.5)$ and $z\mapsto.5+.15(z-.5)$. The same unknown permutation governs every context and the chain. Because the experiments repeat the same type of weak context, the diagnostic-access sweep tests dilution of useful queries rather than learning unknown cross-context correspondences.

The six exact-model conditions factor query allocation from termination:
\begin{center}\small
\begin{tabular}{lll}\toprule
Code name & Query allocation & Stopping target\\\midrule
\texttt{task\_task} & Task-alternative KL maximin & Partial mapping\\
\texttt{full\_task} & Full-alternative KL maximin & Partial mapping\\
\texttt{full\_full} & Full-alternative KL maximin & Full mapping\\
\texttt{task\_full} & Task-alternative KL maximin & Full mapping\\
\texttt{uniform\_task} & Uniform queries & Partial mapping\\
\texttt{entropy\_task} & Full-mapping posterior information gain & Partial mapping\\\bottomrule
\end{tabular}
\end{center}
For entropy allocation, a uniform prior on the finite candidate family is updated by the exact likelihood. Each query's expected reduction of full-mapping entropy is computed exactly; a maximizing query is selected, then mixed with the same $t^{-1/2}$ uniform exploration as the KL methods. The entropy baseline's stopping rule is the same frequentist confidence certificate, not a posterior threshold. All methods use deterministic index-order tie breaking. Each seed draws a new uniform true permutation and paired uniforms for query and outcome sampling. Draws are generated for every replica at every time step, including replicas that have already stopped, preserving paired streams across methods until the last relevant active replica.

Values and regrets are calculated from the known chain, without rollout noise. A cap hit records the current maximum-likelihood policy but labels the run censored and uncertified. Every mean stopping count includes cap hits at the cap. Thus it estimates the restricted mean $\E\min(\tau,\mathrm{cap})$, not an uncensored mean when cap hits occur. Every incorrect-certificate count excludes censored policy guesses.

\begin{center}\small
\begin{tabular}{lrrr}\toprule
Suite & Conditions & Replicas per condition & Total runs\\\midrule
Nuisance sweep & $4\times6$ & 200 & 4,800\\
Confidence sweep & $5\times3$ & 200 & 3,000\\
Estimated kernels & $5\times2$ & $20\times24$ & 4,800\\
Context mismatch & $6\times2$ & 200 & 2,400\\
All-effects control & 4 & 200 & 800\\\midrule
Total &&&15,800\\\bottomrule
\end{tabular}
\end{center}
The exact nuisance sweep uses $\gamma\in\{.01,.02,.04,.08\}$, nine contexts, $\delta=.05$, and cap 50,000. The confidence sweep uses $\gamma=.02$, two contexts, $\delta\in\{.1,.01,.001,10^{-6},10^{-10}\}$, cap 50,000, and task/task, full/task, and entropy/task. The all-effects control uses $\gamma=.04$, two contexts, $R=(0,1,2,3)$, $H=4$, and cap 50,000. Its TDG and full/full runs are exactly identical because the two objectives coincide.

\subsection{Estimated likelihood models}
For a source size $n\in\{256,1024,4096,16384,65536\}$, each anonymous demonstrator supplies $n$ observations for its single fixed latent effect at the one calibration context. With counts $C_{j,y}$ and four outcomes, the estimated distribution is
\begin{equation}
\widehat P_j(y)=\frac{C_{j,y}+1}{n+4}.
\end{equation}
Total passive data are $4n$ transitions per pretrained model. The identity of the demonstrator supplies an effect grouping; it does not supply the executable command. We generate 20 independent pretrained models at each $n$, and 24 calibration replicas per model for both correction modes, at $\gamma=.02$, $\delta=.05$, and cap 3,000. The same model is used to compute both nominal optimal designs and nominal likelihoods. Deployment values remain those of the exact chain.

The correction uses the oracle simulator quantity
\begin{equation}
\eta=\max_{j,y}|\log\widehat P_j(y)-\log P_j(y)|.
\end{equation}
This is a diagnostic study of the robustness theorem with its assumption satisfied, not a practical estimator of $\eta$. No general video encoder, latent clustering objective, or end-to-end dynamics learner is trained. Models and radii are serialized in \texttt{learned\_models.json}. Seed-level intervals resample the 20 pretrained models, not all 480 calibration episodes as independent observations.

\begin{figure}[t]
\centering\includegraphics[width=.94\linewidth]{../figures/learned_robustness.pdf}
\caption{Conservative calibration with estimated categorical effects. Each point includes 20 independently estimated models and 24 calibration replicas per model. Robust radii are computed from simulator truth. The low-data robust rule often reaches the cap without a certificate.}\label{fig:learned}
\end{figure}

\begin{center}\small
\begin{tabular}{rrrrr}\toprule
$n$ per effect & Nominal mean & Robust capped mean & Robust certified & Robust wrong cert.\\\midrule
256 & 63.0 & 3000.0 & 0/480 & 0\\
1024 & 62.9 & 3000.0 & 0/480 & 0\\
4096 & 61.0 & 2077.2 & 173/480 & 0\\
16384 & 60.9 & 342.4 & 461/480 & 0\\
65536 & 61.2 & 90.6 & 480/480 & 0\\\bottomrule
\end{tabular}
\end{center}
The nominal method certifies all runs; it produces one wrong certificate at $n=256$ and none at the other source sizes. These favorable counts on estimated kernels do not establish nominal-model validity. The following constructed shift demonstrates why a guarantee needs more than apparently good empirical calibration.

\subsection{Deliberate semantic mismatch}
Let $S$ exchange latent effects $0$ and $1$. The nominal calibration kernels are
\begin{equation}
\widehat P_j=(1-\alpha)P_j+\alpha P_{S(j)},\qquad
\alpha\in\{0,.1,.3,.5,.7,1\}.
\end{equation}
Actual calibration observations are generated by $P$, and deployment requires the original effect identities. Thus $\alpha=1$ is a consistent alternative permutation for calibration alone but inconsistent with the unchanged deployment semantics. It is not globally equivalent to relabeling the entire model. At $\alpha=.5$, effects $0$ and $1$ are identical in the nominal calibration model, so task identification can fail to terminate. The nominal model's design LP is still solvable with zero worst-case rate; no identifiable-instance guarantee is asserted there.

\begin{figure}[t]
\centering\includegraphics[width=.94\linewidth]{../figures/misspecification.pdf}
\caption{Incorrect certificates under inconsistent calibration and deployment semantics. There are 200 runs per point. The robust rule produces no incorrect certificates, but does so by abstaining at the cap for every nonzero mixture weight. This is a validity demonstration with substantial conservatism, not an efficient repair algorithm.}\label{fig:mismatch}
\end{figure}
At $\alpha=.7$ and $1$, every nominal run certifies an incorrect controller. At $.3$, one of 200 does. At $.5$, no nominal run certifies by the cap. The robust mode computes its log-error radius from the true simulator, produces zero wrong certificates, and certifies only the zero-mismatch condition.

\subsection{Confidence sweep and optimal information rates}
\begin{figure}[t]
\centering\includegraphics[width=.94\linewidth]{../figures/confidence_scaling.pdf}
\caption{Stricter requested confidence favors the task maximin design over entropy reduction in this fixed instance. Each point averages 200 runs. These experiments measure stopping costs; they cannot resolve empirical error probabilities as small as $10^{-10}$.}\label{fig:confidence}
\end{figure}
\begin{center}\small
\begin{tabular}{rrrr}\toprule
$\delta$ & Task/task & Entropy/task & Full/task\\\midrule
$.1$ & 64.4 & 55.4 & 475.9\\
$.01$ & 84.5 & 74.9 & 783.6\\
$.001$ & 103.3 & 96.6 & 1127.2\\
$10^{-6}$ & 158.3 & 170.2 & 2364.6\\
$10^{-10}$ & 224.0 & 375.5 & 4232.1\\\bottomrule
\end{tabular}
\end{center}
We also solve the exact information-design LP over 15 logarithmically spaced $\gamma$ values from $.001$ to $.08$. Task information rates stay bounded away from zero in the small-$\gamma$ regime while full-identification rates vanish quadratically. These are numerical optimization calculations, not additional Monte Carlo replicas. The file \texttt{component\_rates.csv} evaluates the local rates for four context strengths and their cumulative inverse-rate sum. It illustrates the direct-sum formula only; it is not a learned cross-context alignment experiment.

\begin{figure}[t]
\centering\includegraphics[width=.94\linewidth]{../figures/information_geometry.pdf}
\caption{Inverse optimal information rate, computed by the exact alternative LP. This structural calculation explains the nuisance-separation trend independently of any particular stopping-rule overhead.}\label{fig:rates}
\end{figure}

\subsection{Uncertainty estimates, software, and verification}
Summary intervals use 4,000 deterministic percentile bootstrap resamples with seed 92319. Ordinary suites resample calibration replicas; learned-model suites resample pretrained-model clusters. Reported intervals describe uncertainty in the restricted mean, not coverage of the anytime confidence set. The exact raw files also record whether the true mapping ever left the confidence set before stopping. This event can occur more often than a wrong task answer, because distinct mappings can share the same answer.

The implementation uses Python, NumPy, SciPy's HiGHS linear-program solver, and Matplotlib. No GPU or external model API is used. The reproduction scripts record software versions and elapsed time. Tests compare the shortest relevant cycle against brute-force enumeration, cycle and permutation information minima, cutting-plane and enumerated LP objectives, exact negative-control trajectories, confidence thresholds, model-error algebra, and cap handling. Automated tests are numerical and software checks, not formal proof verification.

\end{document}
